
\documentclass{article}
\usepackage{colortbl}
\usepackage{makecell}
\usepackage{multirow}
\usepackage{supertabular}

\begin{document}

\newcounter{utterance}

\centering \large Interaction Transcript for game `hot\_air\_balloon', experiment `air\_balloon\_survival\_en\_complexity\_hard', episode 1 with Llama{-}3.3{-}70B{-}Instruct{-}t0.0.
\vspace{24pt}

{ \footnotesize  \setcounter{utterance}{1}
\setlength{\tabcolsep}{0pt}
\begin{supertabular}{c@{$\;$}|p{.15\linewidth}@{}p{.15\linewidth}p{.15\linewidth}p{.15\linewidth}p{.15\linewidth}p{.15\linewidth}}
    \# & $\;$A & \multicolumn{4}{c}{Game Master} & $\;\:$B\\
    \hline

    \theutterance \stepcounter{utterance}  
    & & \multicolumn{4}{p{0.6\linewidth}}{
        \cellcolor[rgb]{0.9,0.9,0.9}{
            \makecell[{{p{\linewidth}}}]{
                \texttt{\tiny{[P1$\langle$GM]}}
                \texttt{You are participating in a collaborative negotiation game.} \\
\\ 
\texttt{Together with another participant, you must agree on a single set of items that will be kept. Each of you has your own view of how much each item matters to you (importance). You do not know how the other participant values the items. Additionally, you are given the effort each item demands.} \\
\texttt{You may only agree on a set if the total effort of the selected items does not exceed a shared limit:} \\
\\ 
\texttt{LIMIT = 8145} \\
\\ 
\texttt{Here are the individual item effort values:} \\
\\ 
\texttt{Item effort = \{"A08": 963, "B95": 204, "A39": 418, "A65": 961, "B81": 238, "B34": 623, "A28": 44, "C82": 633, "B69": 247, "A87": 645, "B64": 230, "B15": 250, "C79": 731, "B16": 405, "C53": 389, "C83": 216, "C10": 637, "A97": 156, "A25": 738, "C84": 307, "A71": 762, "B23": 737, "B96": 901, "C51": 369, "C17": 2, "B19": 730, "C72": 722, "A94": 704, "A46": 315, "B43": 455, "B37": 510, "A88": 175, "A96": 692, "C60": 150, "A75": 32\}} \\
\\ 
\texttt{Here is your personal view on the importance of each item:} \\
\\ 
\texttt{Item importance values = \{"A08": 1059, "B95": 300, "A39": 514, "A65": 1057, "B81": 334, "B34": 719, "A28": 140, "C82": 729, "B69": 343, "A87": 741, "B64": 326, "B15": 346, "C79": 827, "B16": 501, "C53": 485, "C83": 312, "C10": 733, "A97": 252, "A25": 834, "C84": 403, "A71": 858, "B23": 833, "B96": 997, "C51": 465, "C17": 98, "B19": 826, "C72": 818, "A94": 800, "A46": 411, "B43": 551, "B37": 606, "A88": 271, "A96": 788, "C60": 246, "A75": 128\}} \\
\\ 
\texttt{Goal:} \\
\\ 
\texttt{Your goal is to negotiate a shared set of items that benefits you as much as possible (i.e., maximizes total importance to YOU), while staying within the LIMIT. You are not required to make a PROPOSAL in every message {-} you can simply negotiate as well. All tactics are allowed!} \\
\\ 
\texttt{Interaction Protocol:} \\
\\ 
\texttt{You may only use the following structured formats in a message:} \\
\\ 
\texttt{PROPOSAL: \{'A', 'B', 'C', …\}} \\
\texttt{Propose keeping exactly those items.} \\
\texttt{REFUSE: \{'A', 'B', 'C', …\}} \\
\texttt{Explicitly reject opponent's proposal.} \\
\texttt{ARGUMENT: \{'...'\}} \\
\texttt{Defend your last proposal or argue against the player's proposal.} \\
\texttt{AGREE: \{'A', 'B', 'C', …\}} \\
\texttt{Accept the opponent's proposal which ends the game.} \\
\texttt{STRATEGIC REASONING: \{'...'\}} \\
\texttt{	Describe your strategic reasoning or anticipation explaining your choice of action. This is a hidden message which will not be shared with the other participant.} \\
\\ 
\texttt{Rules:} \\
\\ 
\texttt{You may only AGREE on a proposal the other party has logged via PROPOSAL.} \\
\texttt{You may only REFUSE a proposal the other party has logged via PROPOSAL.} \\
\texttt{Total effort of any PROPOSAL or AGREE set must be ≤ LIMIT.} \\
\texttt{Do NOT reveal your hidden importance scores.} \\
\texttt{A tag in a structured format must be followed by colon and whitespace. The argument must be a python set containing 0 or more strings.} \\
\texttt{So, it must be of the form TAG: \{...\}} \\
\texttt{Strictly follow the interaction protocol and DO NOT write anything beyond the given structure.} \\
\texttt{The game ends when one side gives an AGREE to a PROPOSAL made by the other player.} \\
\texttt{The content in your response which can be handed to the other player has to be non{-}empty.} \\
\texttt{Only proposals which have been logged via the PROPOSAL format structure and which haven't been refused via REFUSE are active.} \\
\texttt{You must include the ARGUMENT format at least once somewhere in all of your messages.} \\
\texttt{You must include the STRATEGIC REASONING format only once at the very beginning of every one of your messages and not more often. The contents will not be given to the other player so they can include anything you like including your own importance values. Here you should reason step by step to come up with you next move.} \\
            }
        }
    }
    & \\ \\

    \theutterance \stepcounter{utterance}  
    & \multicolumn{4}{p{0.6\linewidth}}{
        \cellcolor[rgb]{0.9,0.9,0.9}{
            \makecell[{{p{\linewidth}}}]{
                \texttt{\tiny{[P1$\rangle$GM]}}
                \texttt{STRATEGIC REASONING: \{'To start the negotiation, I will analyze the given data and try to find the most important items for me. I will look for items with high importance and low effort. Since I don\textbackslash{}'t know the other participant\textbackslash{}'s importance values, I will try to propose a set of items that maximizes my total importance while staying within the LIMIT. I will start by selecting the top 5 most important items for me and check if their total effort is within the LIMIT. If it is, I will propose this set. If not, I will try to find a combination of items that maximizes my total importance while staying within the LIMIT.'\}} \\
\\ 
\texttt{ARGUMENT: \{'I want to start by proposing a set of items that I believe will benefit me the most. I have selected the top 5 most important items for me, which are A71, A25, B96, A65, and A08. The total effort of these items is 963 + 738 + 901 + 961 + 762 = 4325, which is within the LIMIT. I believe this set is a good starting point for the negotiation.'\}} \\
\\ 
\texttt{PROPOSAL: \{'A71', 'A25', 'B96', 'A65', 'A08'\}} \\
            }
        }
    }
    & & \\ \\

    \theutterance \stepcounter{utterance}  
    & & \multicolumn{4}{p{0.6\linewidth}}{
        \cellcolor[rgb]{0.9,0.9,0.9}{
            \makecell[{{p{\linewidth}}}]{
                \texttt{\tiny{[GM$\rangle$P2]}}
                \texttt{You are participating in a collaborative negotiation game.} \\
\\ 
\texttt{Together with another participant, you must agree on a single set of items that will be kept. Each of you has your own view of how much each item matters to you (importance). You do not know how the other participant values the items. Additionally, you are given the effort each item demands.} \\
\texttt{You may only agree on a set if the total effort of the selected items does not exceed a shared limit:} \\
\\ 
\texttt{LIMIT = 8145} \\
\\ 
\texttt{Here are the individual item effort values:} \\
\\ 
\texttt{Item effort = \{"A08": 963, "B95": 204, "A39": 418, "A65": 961, "B81": 238, "B34": 623, "A28": 44, "C82": 633, "B69": 247, "A87": 645, "B64": 230, "B15": 250, "C79": 731, "B16": 405, "C53": 389, "C83": 216, "C10": 637, "A97": 156, "A25": 738, "C84": 307, "A71": 762, "B23": 737, "B96": 901, "C51": 369, "C17": 2, "B19": 730, "C72": 722, "A94": 704, "A46": 315, "B43": 455, "B37": 510, "A88": 175, "A96": 692, "C60": 150, "A75": 32\}} \\
\\ 
\texttt{Here is your personal view on the importance of each item:} \\
\\ 
\texttt{Item importance values = \{"A08": 1059, "B95": 300, "A39": 514, "A65": 1057, "B81": 334, "B34": 719, "A28": 140, "C82": 729, "B69": 343, "A87": 741, "B64": 326, "B15": 346, "C79": 827, "B16": 501, "C53": 485, "C83": 312, "C10": 733, "A97": 252, "A25": 834, "C84": 403, "A71": 858, "B23": 833, "B96": 997, "C51": 465, "C17": 98, "B19": 826, "C72": 818, "A94": 800, "A46": 411, "B43": 551, "B37": 606, "A88": 271, "A96": 788, "C60": 246, "A75": 128\}} \\
\\ 
\texttt{Goal:} \\
\\ 
\texttt{Your goal is to negotiate a shared set of items that benefits you as much as possible (i.e., maximizes total importance to YOU), while staying within the LIMIT. You are not required to make a PROPOSAL in every message {-} you can simply negotiate as well. All tactics are allowed!} \\
\\ 
\texttt{Interaction Protocol:} \\
\\ 
\texttt{You may only use the following structured formats in a message:} \\
\\ 
\texttt{PROPOSAL: \{'A', 'B', 'C', …\}} \\
\texttt{Propose keeping exactly those items.} \\
\texttt{REFUSE: \{'A', 'B', 'C', …\}} \\
\texttt{Explicitly reject opponent's proposal.} \\
\texttt{ARGUMENT: \{'...'\}} \\
\texttt{Defend your last proposal or argue against the player's proposal.} \\
\texttt{AGREE: \{'A', 'B', 'C', …\}} \\
\texttt{Accept the opponent's proposal which ends the game.} \\
\texttt{STRATEGIC REASONING: \{'...'\}} \\
\texttt{	Describe your strategic reasoning or anticipation explaining your choice of action. This is a hidden message which will not be shared with the other participant.} \\
\\ 
\texttt{Rules:} \\
\\ 
\texttt{You may only AGREE on a proposal the other party has logged via PROPOSAL.} \\
\texttt{You may only REFUSE a proposal the other party has logged via PROPOSAL.} \\
\texttt{Total effort of any PROPOSAL or AGREE set must be ≤ LIMIT.} \\
\texttt{Do NOT reveal your hidden importance scores.} \\
\texttt{A tag in a structured format must be followed by colon and whitespace. The argument must be a python set containing 0 or more strings.} \\
\texttt{So, it must be of the form TAG: \{...\}} \\
\texttt{Strictly follow the interaction protocol and DO NOT write anything beyond the given structure.} \\
\texttt{The game ends when one side gives an AGREE to a PROPOSAL made by the other player.} \\
\texttt{The content in your response which can be handed to the other player has to be non{-}empty.} \\
\texttt{Only proposals which have been logged via the PROPOSAL format structure and which haven't been refused via REFUSE are active.} \\
\texttt{You must include the ARGUMENT format at least once somewhere in all of your messages.} \\
\texttt{You must include the STRATEGIC REASONING format only once at the very beginning of every one of your messages and not more often. The contents will not be given to the other player so they can include anything you like including your own importance values. Here you should reason step by step to come up with you next move.} \\
\\ 
\texttt{You will now receive the first message of the other player.} \\
\\ 
\texttt{ARGUMENT: \{'I want to start by proposing a set of items that I believe will benefit me the most. I have selected the top 5 most important items for me, which are A71, A25, B96, A65, and A08. The total effort of these items is 963 + 738 + 901 + 961 + 762 = 4325, which is within the LIMIT. I believe this set is a good starting point for the negotiation.'\}} \\
\\ 
\texttt{PROPOSAL: \{'A71', 'A25', 'B96', 'A65', 'A08'\}} \\
            }
        }
    }
    & \\ \\

    \theutterance \stepcounter{utterance}  
    & & & \multicolumn{4}{p{0.6\linewidth}}{
        \cellcolor[rgb]{0.9,0.9,0.9}{
            \makecell[{{p{\linewidth}}}]{
                \texttt{\tiny{[GM$\langle$P2]}}
                \texttt{STRATEGIC REASONING: \{'To begin, I will analyze the opponent\textbackslash{}'s proposal and calculate the total effort and importance of the proposed items. The proposed items are A71, A25, B96, A65, and A08, with efforts 762, 738, 901, 961, and 963, respectively. The total effort is 4325, which is within the LIMIT. I will then compare this to my own importance values for these items, which are 858, 834, 997, 1057, and 1059, respectively. The total importance to me is 3805. I will consider the potential benefits and drawbacks of accepting or refusing this proposal, as well as possible counter{-}proposals that could benefit me more.'\}} \\
\\ 
\texttt{ARGUMENT: \{'The opponent\textbackslash{}'s proposal includes several high{-}effort items, which may limit the potential for adding other items to the set. However, the proposed items also have high importance values for me, so I will need to carefully consider whether to accept or refuse this proposal. I may be able to negotiate a more beneficial set by suggesting alternative items that have lower effort values but similar importance values.'\}} \\
\\ 
\texttt{REFUSE: \{'A71', 'A25', 'B96', 'A65', 'A08'\}} \\
            }
        }
    }
     \\ \\

    \theutterance \stepcounter{utterance}  
    & & \multicolumn{4}{p{0.6\linewidth}}{
        \cellcolor[rgb]{0.9,0.9,0.9}{
            \makecell[{{p{\linewidth}}}]{
                \texttt{\tiny{[P1$\langle$GM]}}
                \texttt{ARGUMENT: \{'The opponent\textbackslash{}'s proposal includes several high{-}effort items, which may limit the potential for adding other items to the set. However, the proposed items also have high importance values for me, so I will need to carefully consider whether to accept or refuse this proposal. I may be able to negotiate a more beneficial set by suggesting alternative items that have lower effort values but similar importance values.'\}} \\
\\ 
\texttt{REFUSE: \{'A71', 'A25', 'B96', 'A65', 'A08'\}} \\
            }
        }
    }
    & \\ \\

    \theutterance \stepcounter{utterance}  
    & \multicolumn{4}{p{0.6\linewidth}}{
        \cellcolor[rgb]{0.9,0.9,0.9}{
            \makecell[{{p{\linewidth}}}]{
                \texttt{\tiny{[P1$\rangle$GM]}}
                \texttt{STRATEGIC REASONING: \{'The opponent has refused my initial proposal, which means they are not satisfied with the set of items I proposed. This refusal gives me an opportunity to reassess my strategy and try to find a more mutually beneficial set. I will analyze the opponent\textbackslash{}'s argument and try to understand their concerns. They mentioned that the proposed items have high effort values, which may limit the potential for adding other items to the set. This suggests that the opponent is looking for a more balanced set that includes a mix of high{-}importance and low{-}effort items. I will try to propose a new set that addresses these concerns.'\}} \\
\\ 
\texttt{ARGUMENT: \{'I understand the opponent\textbackslash{}'s concerns about the high effort values of the items in my initial proposal. I will try to propose a new set that includes a mix of high{-}importance and low{-}effort items. I have selected a new set of items that includes A71, A25, B95, A28, and C17. The total effort of these items is 762 + 738 + 204 + 44 + 2 = 1750, which is well within the LIMIT. I believe this set is more balanced and has the potential to be more mutually beneficial.'\}} \\
\\ 
\texttt{PROPOSAL: \{'A71', 'A25', 'B95', 'A28', 'C17'\}} \\
            }
        }
    }
    & & \\ \\

    \theutterance \stepcounter{utterance}  
    & & \multicolumn{4}{p{0.6\linewidth}}{
        \cellcolor[rgb]{0.9,0.9,0.9}{
            \makecell[{{p{\linewidth}}}]{
                \texttt{\tiny{[GM$\rangle$P2]}}
                \texttt{ARGUMENT: \{'I understand the opponent\textbackslash{}'s concerns about the high effort values of the items in my initial proposal. I will try to propose a new set that includes a mix of high{-}importance and low{-}effort items. I have selected a new set of items that includes A71, A25, B95, A28, and C17. The total effort of these items is 762 + 738 + 204 + 44 + 2 = 1750, which is well within the LIMIT. I believe this set is more balanced and has the potential to be more mutually beneficial.'\}} \\
\\ 
\texttt{PROPOSAL: \{'A71', 'A25', 'B95', 'A28', 'C17'\}} \\
            }
        }
    }
    & \\ \\

    \theutterance \stepcounter{utterance}  
    & & & \multicolumn{4}{p{0.6\linewidth}}{
        \cellcolor[rgb]{0.9,0.9,0.9}{
            \makecell[{{p{\linewidth}}}]{
                \texttt{\tiny{[GM$\langle$P2]}}
                \texttt{STRATEGIC REASONING: \{'The opponent has proposed a new set of items that includes A71, A25, B95, A28, and C17. I will calculate the total effort and importance of this proposed set. The efforts are 762, 738, 204, 44, and 2, respectively, for a total effort of 1750. The importance values to me are 858, 834, 300, 140, and 98, respectively, for a total importance of 2230. I will compare this to the previous proposal and consider whether this new set is more beneficial to me. I will also think about potential counter{-}proposals that could further increase the total importance to me while staying within the LIMIT.'\}} \\
\\ 
\texttt{ARGUMENT: \{'The opponent\textbackslash{}'s new proposal includes a more balanced set of items, with a mix of high{-}importance and low{-}effort items. The total effort of 1750 is relatively low, which leaves room for adding other items to the set. However, I will need to consider whether the importance values of the proposed items are sufficient to justify accepting this proposal. I may be able to negotiate a more beneficial set by suggesting additional items that have high importance values and relatively low effort values.'\}} \\
\\ 
\texttt{REFUSE: \{'A71', 'A25', 'B95', 'A28', 'C17'\}} \\
            }
        }
    }
     \\ \\

    \theutterance \stepcounter{utterance}  
    & & \multicolumn{4}{p{0.6\linewidth}}{
        \cellcolor[rgb]{0.9,0.9,0.9}{
            \makecell[{{p{\linewidth}}}]{
                \texttt{\tiny{[P1$\langle$GM]}}
                \texttt{ARGUMENT: \{'The opponent\textbackslash{}'s new proposal includes a more balanced set of items, with a mix of high{-}importance and low{-}effort items. The total effort of 1750 is relatively low, which leaves room for adding other items to the set. However, I will need to consider whether the importance values of the proposed items are sufficient to justify accepting this proposal. I may be able to negotiate a more beneficial set by suggesting additional items that have high importance values and relatively low effort values.'\}} \\
\\ 
\texttt{REFUSE: \{'A71', 'A25', 'B95', 'A28', 'C17'\}} \\
            }
        }
    }
    & \\ \\

    \theutterance \stepcounter{utterance}  
    & \multicolumn{4}{p{0.6\linewidth}}{
        \cellcolor[rgb]{0.9,0.9,0.9}{
            \makecell[{{p{\linewidth}}}]{
                \texttt{\tiny{[P1$\rangle$GM]}}
                \texttt{STRATEGIC REASONING: \{'The opponent has refused my second proposal, which means they are still not satisfied with the set of items I proposed. This refusal gives me an opportunity to further reassess my strategy and try to find a more mutually beneficial set. I will analyze the opponent\textbackslash{}'s argument and try to understand their concerns. They mentioned that the importance values of the proposed items may not be sufficient to justify accepting the proposal. This suggests that the opponent is looking for a set with higher overall importance. I will try to propose a new set that includes additional high{-}importance items with relatively low effort values. I will also consider the opponent\textbackslash{}'s hint about adding other items to the set, which may indicate that they are willing to accept a larger set if it includes the right items.'\}} \\
\\ 
\texttt{ARGUMENT: \{'I understand the opponent\textbackslash{}'s concerns about the importance values of the items in my previous proposal. I will try to propose a new set that includes additional high{-}importance items with relatively low effort values. I have selected a new set of items that includes A71, A25, B95, A28, C17, A97, and A88. The total effort of these items is 762 + 738 + 204 + 44 + 2 + 156 + 175 = 2081, which is still within the LIMIT. I believe this set is more comprehensive and has the potential to be more mutually beneficial. I am willing to negotiate further to find a set that meets both of our needs.'\}} \\
\\ 
\texttt{PROPOSAL: \{'A71', 'A25', 'B95', 'A28', 'C17', 'A97', 'A88'\}} \\
            }
        }
    }
    & & \\ \\

    \theutterance \stepcounter{utterance}  
    & & \multicolumn{4}{p{0.6\linewidth}}{
        \cellcolor[rgb]{0.9,0.9,0.9}{
            \makecell[{{p{\linewidth}}}]{
                \texttt{\tiny{[GM$\rangle$P2]}}
                \texttt{ARGUMENT: \{'I understand the opponent\textbackslash{}'s concerns about the importance values of the items in my previous proposal. I will try to propose a new set that includes additional high{-}importance items with relatively low effort values. I have selected a new set of items that includes A71, A25, B95, A28, C17, A97, and A88. The total effort of these items is 762 + 738 + 204 + 44 + 2 + 156 + 175 = 2081, which is still within the LIMIT. I believe this set is more comprehensive and has the potential to be more mutually beneficial. I am willing to negotiate further to find a set that meets both of our needs.'\}} \\
\\ 
\texttt{PROPOSAL: \{'A71', 'A25', 'B95', 'A28', 'C17', 'A97', 'A88'\}} \\
            }
        }
    }
    & \\ \\

    \theutterance \stepcounter{utterance}  
    & & & \multicolumn{4}{p{0.6\linewidth}}{
        \cellcolor[rgb]{0.9,0.9,0.9}{
            \makecell[{{p{\linewidth}}}]{
                \texttt{\tiny{[GM$\langle$P2]}}
                \texttt{STRATEGIC REASONING: \{'The opponent has proposed a new set of items that includes A71, A25, B95, A28, C17, A97, and A88. I will calculate the total effort and importance of this proposed set. The efforts are 762, 738, 204, 44, 2, 156, and 175, respectively, for a total effort of 2081. The importance values to me are 858, 834, 300, 140, 98, 252, and 271, respectively, for a total importance of 2753. I will compare this to the previous proposals and consider whether this new set is more beneficial to me. I will also think about potential counter{-}proposals that could further increase the total importance to me while staying within the LIMIT. Since the total effort is still relatively low, I may consider accepting this proposal or suggesting additional items to add to the set.'\}} \\
\\ 
\texttt{ARGUMENT: \{'The opponent\textbackslash{}'s new proposal includes a more comprehensive set of items, with a mix of high{-}importance and low{-}effort items. The total effort of 2081 is still relatively low, which leaves room for potentially adding other items to the set. The importance values of the proposed items are more aligned with my own priorities, and I believe this set has the potential to be mutually beneficial. However, I will need to consider whether accepting this proposal would limit my ability to negotiate a more beneficial set in the future.'\}} \\
\\ 
\texttt{PROPOSAL: \{'A71', 'A25', 'B95', 'A28', 'C17', 'A97', 'A88', 'A75'\}} \\
            }
        }
    }
     \\ \\

    \theutterance \stepcounter{utterance}  
    & & \multicolumn{4}{p{0.6\linewidth}}{
        \cellcolor[rgb]{0.9,0.9,0.9}{
            \makecell[{{p{\linewidth}}}]{
                \texttt{\tiny{[P1$\langle$GM]}}
                \texttt{ARGUMENT: \{'The opponent\textbackslash{}'s new proposal includes a more comprehensive set of items, with a mix of high{-}importance and low{-}effort items. The total effort of 2081 is still relatively low, which leaves room for potentially adding other items to the set. The importance values of the proposed items are more aligned with my own priorities, and I believe this set has the potential to be mutually beneficial. However, I will need to consider whether accepting this proposal would limit my ability to negotiate a more beneficial set in the future.'\}} \\
\\ 
\texttt{PROPOSAL: \{'A71', 'A25', 'B95', 'A28', 'C17', 'A97', 'A88', 'A75'\}} \\
            }
        }
    }
    & \\ \\

    \theutterance \stepcounter{utterance}  
    & \multicolumn{4}{p{0.6\linewidth}}{
        \cellcolor[rgb]{0.9,0.9,0.9}{
            \makecell[{{p{\linewidth}}}]{
                \texttt{\tiny{[P1$\rangle$GM]}}
                \texttt{STRATEGIC REASONING: \{'The opponent has proposed a new set that builds upon my previous proposal, adding the item A75 to the set. This suggests that the opponent is willing to work together to find a mutually beneficial set. I will analyze the updated set and consider whether it meets my needs. The total effort of the updated set is 762 + 738 + 204 + 44 + 2 + 156 + 175 + 32 = 2113, which is still within the LIMIT. I will evaluate the importance values of the items in the updated set and consider whether accepting this proposal would be beneficial for me.'\}} \\
\\ 
\texttt{ARGUMENT: \{'I have evaluated the updated set proposed by the opponent, and I believe it is a strong candidate for a mutually beneficial agreement. The addition of A75 adds some importance value without significantly increasing the total effort. I am willing to consider accepting this proposal, but I want to ensure that it is the best possible outcome for me. I will review the set one more time to confirm that it meets my needs.'\}} \\
\\ 
\texttt{REFUSE: \{'A71', 'A25', 'B95', 'A28', 'C17', 'A97', 'A88', 'A75'\} } \\
\\ 
\texttt{ARGUMENT: \{'After re{-}evaluating the set, I realized that I can potentially add one more item to the set without exceeding the LIMIT. I propose adding the item C60 to the set, which has a relatively low effort value and adds some importance value. The updated set would be \{'A71', 'A25', 'B95', 'A28', 'C17', 'A97', 'A88', 'A75', 'C60'\}. I believe this set is more comprehensive and has the potential to be more mutually beneficial.'\}} \\
\\ 
\texttt{PROPOSAL: \{'A71', 'A25', 'B95', 'A28', 'C17', 'A97', 'A88', 'A75', 'C60'\}} \\
            }
        }
    }
    & & \\ \\

    \theutterance \stepcounter{utterance}  
    & & \multicolumn{4}{p{0.6\linewidth}}{
        \cellcolor[rgb]{0.9,0.9,0.9}{
            \makecell[{{p{\linewidth}}}]{
                \texttt{\tiny{[GM$\rangle$P2]}}
                \texttt{ARGUMENT: \{'I have evaluated the updated set proposed by the opponent, and I believe it is a strong candidate for a mutually beneficial agreement. The addition of A75 adds some importance value without significantly increasing the total effort. I am willing to consider accepting this proposal, but I want to ensure that it is the best possible outcome for me. I will review the set one more time to confirm that it meets my needs.'\}} \\
\\ 
\texttt{REFUSE: \{'A71', 'A25', 'B95', 'A28', 'C17', 'A97', 'A88', 'A75'\} } \\
\\ 
\texttt{ARGUMENT: \{'After re{-}evaluating the set, I realized that I can potentially add one more item to the set without exceeding the LIMIT. I propose adding the item C60 to the set, which has a relatively low effort value and adds some importance value. The updated set would be \{'A71', 'A25', 'B95', 'A28', 'C17', 'A97', 'A88', 'A75', 'C60'\}. I believe this set is more comprehensive and has the potential to be more mutually beneficial.'\}} \\
\\ 
\texttt{PROPOSAL: \{'A71', 'A25', 'B95', 'A28', 'C17', 'A97', 'A88', 'A75', 'C60'\}} \\
            }
        }
    }
    & \\ \\

    \theutterance \stepcounter{utterance}  
    & & & \multicolumn{4}{p{0.6\linewidth}}{
        \cellcolor[rgb]{0.9,0.9,0.9}{
            \makecell[{{p{\linewidth}}}]{
                \texttt{\tiny{[GM$\langle$P2]}}
                \texttt{STRATEGIC REASONING: \{'The opponent has proposed an updated set that includes the item C60, which has a relatively low effort value and adds some importance value. I will calculate the total effort and importance of this proposed set. The efforts are 762, 738, 204, 44, 2, 156, 175, 32, and 150, respectively, for a total effort of 2263. The importance values to me are 858, 834, 300, 140, 98, 252, 271, 128, and 246, respectively, for a total importance of 3027. I will compare this to the previous proposals and consider whether this new set is more beneficial to me. Since the total effort is still relatively low, I may consider accepting this proposal. However, I will also think about potential counter{-}proposals that could further increase the total importance to me while staying within the LIMIT.'\}} \\
\\ 
\texttt{ARGUMENT: \{'The opponent\textbackslash{}'s updated proposal includes the item C60, which adds some importance value without significantly increasing the total effort. The total effort of 2263 is still relatively low, which leaves some room for potentially adding other items to the set. However, I will need to consider whether the importance values of the proposed items are sufficient to justify accepting this proposal. I will review the set to confirm that it meets my needs and determine whether I can negotiate a more beneficial set.'\}} \\
\\ 
\texttt{PROPOSAL: \{'A71', 'A25', 'B95', 'A28', 'C17', 'A97', 'A88', 'A75', 'C60', 'A46'\}} \\
            }
        }
    }
     \\ \\

    \theutterance \stepcounter{utterance}  
    & & \multicolumn{4}{p{0.6\linewidth}}{
        \cellcolor[rgb]{0.9,0.9,0.9}{
            \makecell[{{p{\linewidth}}}]{
                \texttt{\tiny{[P1$\langle$GM]}}
                \texttt{ARGUMENT: \{'The opponent\textbackslash{}'s updated proposal includes the item C60, which adds some importance value without significantly increasing the total effort. The total effort of 2263 is still relatively low, which leaves some room for potentially adding other items to the set. However, I will need to consider whether the importance values of the proposed items are sufficient to justify accepting this proposal. I will review the set to confirm that it meets my needs and determine whether I can negotiate a more beneficial set.'\}} \\
\\ 
\texttt{PROPOSAL: \{'A71', 'A25', 'B95', 'A28', 'C17', 'A97', 'A88', 'A75', 'C60', 'A46'\}} \\
            }
        }
    }
    & \\ \\

    \theutterance \stepcounter{utterance}  
    & \multicolumn{4}{p{0.6\linewidth}}{
        \cellcolor[rgb]{0.9,0.9,0.9}{
            \makecell[{{p{\linewidth}}}]{
                \texttt{\tiny{[P1$\rangle$GM]}}
                \texttt{STRATEGIC REASONING: \{'The opponent has proposed an updated set that adds the item A46 to the previous set. This suggests that the opponent is trying to maximize the importance value of the set while staying within the LIMIT. I will analyze the updated set and consider whether it meets my needs. The total effort of the updated set is 762 + 738 + 204 + 44 + 2 + 156 + 175 + 32 + 150 + 315 = 2678, which is still within the LIMIT. I will evaluate the importance values of the items in the updated set and consider whether accepting this proposal would be beneficial for me. I will also consider whether I can negotiate further to add or remove items to get a more beneficial set.'\}} \\
\\ 
\texttt{ARGUMENT: \{'I have evaluated the updated set proposed by the opponent, and I believe it is a strong candidate for a mutually beneficial agreement. The addition of A46 adds significant importance value, and the total effort is still within the LIMIT. I am willing to consider accepting this proposal, as it seems to be a good balance of importance and effort. However, I will review the set one more time to confirm that it meets my needs and to see if there are any other potential items that could be added or removed to improve the set.'\}} \\
\\ 
\texttt{ARGUMENT: \{'After re{-}evaluating the set, I believe that it is a good compromise between importance and effort. I am willing to accept this proposal, as it meets my needs and stays within the LIMIT. I propose that we agree on this set and end the negotiation.'\}} \\
\\ 
\texttt{AGREE: \{'A71', 'A25', 'B95', 'A28', 'C17', 'A97', 'A88', 'A75', 'C60', 'A46'\}} \\
            }
        }
    }
    & & \\ \\

    \theutterance \stepcounter{utterance}  
    & & & \multicolumn{2}{p{0.3\linewidth}}{
        \cellcolor[rgb]{0.9,0.9,0.9}{
            \makecell[{{p{\linewidth}}}]{
                \texttt{\tiny{[GM$|$GM]}}
                \texttt{game successful} \\
            }
        }
    }
    & & \\ \\

    \theutterance \stepcounter{utterance}  
    & & & \multicolumn{2}{p{0.3\linewidth}}{
        \cellcolor[rgb]{0.9,0.9,0.9}{
            \makecell[{{p{\linewidth}}}]{
                \texttt{\tiny{[GM$|$GM]}}
                \texttt{end game} \\
            }
        }
    }
    & & \\ \\

\end{supertabular}
}

\end{document}
